% Self-contained editable report. Compile twice with:
% pdflatex -interaction=nonstopmode -halt-on-error CS6886_Report.tex
% No external figures, bibliography program, or custom class is required.
\documentclass[10pt,a4paper]{article}
\usepackage[T1]{fontenc}
\usepackage{lmodern,microtype}
\DisableLigatures{encoding = *, family = tt*}
\usepackage[margin=16mm,top=13mm,bottom=13mm]{geometry}
\usepackage{amsmath,amssymb,booktabs,tabularx,array,multicol}
\usepackage[table]{xcolor}
\usepackage{enumitem,titlesec,fancyhdr,lastpage,xurl}
\usepackage{tikz}
\usetikzlibrary{arrows.meta,positioning}
\usepackage{tcolorbox,listings}
\lstset{basicstyle=\footnotesize\ttfamily,columns=fullflexible,keepspaces=true,
  breaklines=true,showstringspaces=false,backgroundcolor=\color{pale},
  frame=single,rulecolor=\color{accent!35},framerule=.3pt,
  aboveskip=3pt,belowskip=3pt,xleftmargin=3pt,xrightmargin=3pt,
  framexleftmargin=3pt,framexrightmargin=3pt}
\titleformat{\subsection}{\normalsize\bfseries\color{accent}}{\thesubsection}{.5em}{}
\titlespacing*{\subsection}{0pt}{4pt}{2pt}
\usepackage[colorlinks=true,linkcolor=teal!65!black,citecolor=teal!65!black,urlcolor=teal!65!black]{hyperref}
\hypersetup{pdftitle={Compiler Optimizations for Attention and Functional KV Caches},pdfauthor={Nishanth Senthil Kumar}}
\definecolor{accent}{HTML}{176B70}
\definecolor{pale}{HTML}{EFF6F6}
\definecolor{ink}{HTML}{192A35}
\color{ink}
\setlength{\parindent}{0pt}
\setlength{\parskip}{3pt}
\setlength{\abovedisplayskip}{4pt}
\setlength{\belowdisplayskip}{4pt}
\setlength{\abovedisplayshortskip}{3pt}
\setlength{\belowdisplayshortskip}{3pt}
\setlength{\tabcolsep}{4pt}
\renewcommand{\arraystretch}{1.04}
\newcolumntype{Y}{>{\raggedright\arraybackslash}X}
\newcolumntype{P}[1]{>{\raggedright\arraybackslash}p{#1}}
\setlist[itemize]{leftmargin=13pt,itemsep=1.5pt,topsep=2pt,parsep=0pt}
\titleformat{\section}{\large\bfseries\color{accent}}{\thesection}{0.5em}{}
\titlespacing*{\section}{0pt}{4pt}{2pt}
\pagestyle{fancy}
\fancyhf{}
\fancyhead[L]{\footnotesize CS6886 \;|\; Attention compiler optimization}
\fancyhead[R]{\footnotesize 1 November 2026}
\fancyfoot[R]{\footnotesize\thepage\,/\,\pageref{LastPage}}
\renewcommand{\headrulewidth}{0.25pt}
\setlength{\headheight}{13pt}
\newcommand{\code}[1]{\texttt{#1}}
\newcommand{\sel}{\operatorname{select}}
\newcommand{\cat}{\operatorname{concat}}
\newcommand{\up}{\operatorname{update}}
\newcommand{\smallnote}[1]{{\footnotesize #1\par}}
\newtcolorbox{resultbox}{colback=pale,colframe=accent,boxrule=0.45pt,arc=1mm,left=5pt,right=5pt,top=3pt,bottom=3pt}
% Edit/add team members here as appropriate.
\newcommand{\projectauthors}{Nishanth Senthil Kumar (EE23B049) \quad IIT Madras}


\begin{document}
\thispagestyle{plain}
{\LARGE\bfseries Compiler Optimizations for Attention\\[2pt]and Functional KV Caches\par}
{\small CS6886: Systems Engineering for Deep Learning \hfill 1 November 2026\par}
{\small\projectauthors\par}


\section{Why the two IREE flags remove the original GQA repeat}
\textbf{Starting graph.} GQA shares $H_{kv}$ key/value heads among $H_q$ query heads~\cite{gqa,mqa}. For 8 query heads and 2 KV heads, $g=4$: query heads 0--3 use KV head 0 and heads 4--7 use KV head 1. With batch omitted, $Q:[8,M,D]$ and $K,V:[2,N,D]$:
\begin{lstlisting}
# Natural decomposed attention: expanded tensors have 8 heads.
Kr = repeat_interleave(K, 4, dim=0)   # [8, N, D]
Vr = repeat_interleave(V, 4, dim=0)
P  = softmax((Q @ Kr.transpose(-2, -1)) / sqrt(D), dim=-1)
O  = P @ Vr
\end{lstlisting}
The saved Linalg IR first broadcasts K/V to $[2,4,N,D]$, collapses the two head axes to 8, then feeds named \code{linalg.batch\_matmul} operations. The broadcast is a tensor-producing operation; reshaping alone is not evidence of a copy.

\subsection*{Flag 1: expose matmul as a generic contraction}
{\footnotesize\ttfamily --iree-opt-generalize-matmul\par}
The \code{GeneralizeLinalgNamedOps} pass rewrites named matmuls as \code{linalg.generic}, exposing parallel/reduction loops, operand indexing maps and multiply-add arithmetic~\cite{ireesource,ireedocs}. A reduction sums many feature products into one output:
\begin{lstlisting}
# Schematic loop meaning of linalg.batch_matmul(Q, Kr.T):
for h, m, n:                    # parallel output indices
    S[h,m,n] = sum_d(Q[h,m,d] * Kr[h,n,d])
\end{lstlisting}
\textbf{Why needed:} the reshape pass has rules for generic operations; the pinned implementation excludes propagation through named Linalg operations. Generalization exposes the structure those rules need. It does not itself guarantee fusion or a faster kernel.

\subsection*{Flag 2: propagate reshapes through reductions}
{\footnotesize\ttfamily --iree-dispatch-creation-enable-aggressive-reshape-movement\par}
The reduction-crossing option in \code{BubbleUpExpandShapes} permits expand/collapse propagation through eligible generic reductions, rather than restricting it to all-parallel operations~\cite{ireesource,ireedocs}.

\textbf{Effect here:} the flattened head axis can be split into \emph{KV head}$\times$\emph{sharing group}. Existing elementwise fusion can then combine the broadcast with the contraction and compose its indexing map. The resulting loop meaning is:
\begin{lstlisting}
for hk in range(2):
  for j in range(4):
    h = 4*hk + j
    S[h,m,n] = sum_d(Q[h,m,d] * K[hk,n,d])  # no Kr
# The probability-times-V contraction analogously reads V[hk].
\end{lstlisting}
This equals $h_{kv}=\lfloor h_q/4\rfloor$, but the compiler can retain separate \code{hk,j} loops without emitting an explicit integer divide. This changes loop shapes/indexing while preserving the complete dot product. The documentation is therefore under general Linalg/shape-fusion passes, rather than a GQA-specific rule.

\subsection*{Evidence and implication for our original pass}
{\small\begin{tabularx}{\linewidth}{@{}YYrr@{}}\toprule
Generalize matmul & Aggressive reshape & Broadcast calls & All dispatch calls\\\midrule
Off&Off&2&5\\On&Off&2&5\\Off&On&2&5\\On&On&0&3\\\bottomrule\end{tabularx}}
\smallnote{Fresh O3 CPU compile-to-Flow check: $H_q=8,H_{kv}=2$. Observed order: generalization, reshape movement, then elementwise fusion. Final contractions consume compact $[2,N,D]$ K/V. Earlier ablation: 20 fixtures, 160 compilations and 96 numerical checks.}

\textbf{What becomes redundant?} Our pass manually constructs grouped contractions; these flags plus fusion already eliminate the tested repeats. The original 1.86$\times$ default CPU gain therefore does not show an additional capability over strong IREE. O3 alone is insufficient. This is not all-graph coverage: MQA can need fewer flags, and extra or escaping users can retain expanded tensors.

\newpage
\section{Functional-cache transformations: what changes and why}
\smallnote{Examples are explanatory pseudocode. \code{update(C,p,U)} returns a functional tensor value with U replacing rows $p:p+t$; it does not mean every backend must allocate/copy the whole C. Let $C:[P,H,D]$ and $U:[t,H,D]$.}

\subsection*{2.1 Conditional read forwarding: avoid building an unused full version}
\textbf{Before:} each candidate forms an updated cache, but downstream sparse attention only gathers a few rows.
\begin{lstlisting}
T = update(C, p, U)             # logical full updated version
R = gather(T, ids)
# After: conditional loads from original history or the patch.
R = stack([U[i-p] if p <= i < p+t else C[i] for i in ids])
\end{lstlisting}
Example: $P=10000,p=9000,t=1$, \code{ids=[0,9000,9999]} returns \code{[C[0],U[0],C[9999]]}. If T has only supported read consumers, its full materialization can disappear. \textbf{Use case:} several branches share a large pool but each reads a small subset. \textbf{Limits:} required full-cache outputs remain; supported types, indices, slices and uses are checked. Read-over-write is classical array algebra~\cite{mccarthy,arrays}, not a new identity.

\subsection*{2.2 Masked gather forwarding: a different implementation of the same law}
Instead of per-element conditional control flow, gather a compact base and patch and select their elements. The implemented masked variant uses a safe patch index even outside the update:
\begin{lstlisting}
inside = (p <= ids) & (ids < p+t)
safe   = where(inside, ids-p, 0)  # U must have at least 1 row
base   = gather(C, ids)
patch  = gather(U, safe)
R      = where(expand_mask(inside), patch, base)
\end{lstlisting}
\textbf{Why the safe index matters:} a select may evaluate both operands; an out-of-bounds patch read cannot be justified by a later mask. \textbf{Tradeoff:} masks, two gathers and selection add work. A demand/capacity cost gate rejects dense cases; the masked form also has an irregular-shape CUDA compile failure with a conditional fallback. It is not universally faster.

\subsection*{2.3 Winner-cache factoring: select the small update before commit}
Two candidates share the same history but propose different new rows. The final choice is a scalar Boolean:
\begin{lstlisting}
# Before: selection consumes two full appended cache values.
Ta, Tb = concat(C, a), concat(C, b)
committed = select(winner_is_a, Ta, Tb)
# After: move the same selection to the small appended rows.
committed = concat(C, select(winner_is_a, a, b))
\end{lstlisting}
Example: with 4096 cached rows and one new row, selection moves from 4097-row tensors to one-row tensors. The pass also supports equal-region insert-slice forms. It checks a common base, axis/region, types and scalar condition; different histories or elementwise conditions reject. \textbf{Use case:} branching inference commit. The pretrained driver has 60 concat matches; ordinary decode has zero. Both candidate attention computations still need their inputs, so this does not prove all candidate caches or copies disappear.

\subsection*{2.4 Adjacent-update coalescing: build one small patch}
\begin{lstlisting}
# Before: intermediate versions are not observed elsewhere.
C1 = update(C, 100, a)           # a has 1 row
C2 = update(C1, 101, b)          # b has 1 row
# After:
C2 = update(C, 100, concat(a, b))
\end{lstlisting}
\textbf{Use case:} supported multi-token update chains before forwarding. Adjacent regions become one small patch; fewer intermediate updates can expose later elimination. Guards require compatible slices, positive sizes, unit strides, no gap/overlap, a single-use intermediate and a sufficiently small patch. The implementation targets copy-prone shared bases; an already-in-place single-use update may not benefit. Earlier CPU chains improve 2.17--2.34$\times$; this is not an established general GPU gain.

\subsection*{2.5 Gather-index normalization: preserve the right batch coordinates}
For source $C:[P,H,D]$ and indices $[F,N,1]$, the last1 is an index-component axis, not another batch axis. The required read is $R[f,n,h,d]=C[\mathrm{ids}[f,n,0],h,d]$.
\begin{lstlisting}
ids2 = reshape(ids, [F, N])      # remove ONLY the unit component
R    = gather(C, ids2)           # preserve f and n coordinates
\end{lstlisting}
This normalizes a supported depth-one gather before vectorization, avoiding the diagnosed index-projection defect. It checks the gather dimension map and matching batch/output shapes. It is a correctness repair, not head sharing; arbitrary index depth/layouts are outside this rewrite.

\newpage
\section{Ownership, projection and numerical transformations}

\subsection*{3.1 Snapshot protection and borrowed-buffer copy-on-write}
\begin{lstlisting}
# Tensor contract: return the OLD C alongside an updated-cache result.
T = update(C, p, U)
return attention(T), C
# Required ownership if bufferized T would alias borrowed C:
T_buf = clone(C_buf); T_buf[p:p+t] = U
return attention(T_buf), C_buf
\end{lstlisting}
An in-place update through the same buffer would corrupt the returned snapshot. The early tensor pass preserves the old observation using an optimization barrier; the phase-local Stream repair follows casts to non-consuming imports and clones the update target in the supported ownership pattern. The code above illustrates that ownership requirement, not a promise to clone every update. If forwarding makes T unnecessary, a full copy can be avoided. Incorrect stock outputs are excluded from speedup baselines. Cache sharing/COW has earlier systems precedents~\cite{paged,forkkv}.

\subsection*{3.2 Latent projection placement: project the query instead of every key}
Let $C:[N,r]$ be an already-compressed cache, $W_K:[r,d]$, and $q:[1,d]$.
\begin{lstlisting}
# Before: project all N cached vectors, then compare with q.
K = C @ Wk                      # [N, d]
scores = q @ K.T
# After: one query projection and a latent-width comparison.
scores = (q @ Wk.T) @ C.T        # [1, r] @ [r, N]
\end{lstlisting}
For $N=10000,r=64,d=512$, this avoids forming 10000 full-width keys and replaces their up-projections with one query projection. The value-side identity is $A(CW_V)=(AC)W_V$. \textbf{Use case:} latent-attention decode, not arbitrary standard KV caches. The pass uses shape/use and cost checks and an explicit reassociation contract: real-number equivalence does not imply bitwise FP equality. It can regress for long queries. MLA absorption is established prior work~\cite{deepseek}.

\subsection*{3.3 Partial-RoPE factoring: absorb content while retaining rotation}
Rotary position transforms generally cannot commute through $W_K$. For an explicit content/rotary split, preserve the rotary contribution and transform only the content term~\cite{deepseek,rope}:
\begin{lstlisting}
# qC:512 content dims, qR:32 rotary dims; C:64 latent dims.
s = (qC @ Wk.T) @ C.T + qR @ kR.T
P = softmax(s / sqrt(512 + 32))  # preserve ORIGINAL scale
\end{lstlisting}
This equals $[q_C,q_R][CW_K,k_R]^T$. The rewritten content width 64 does not justify changing the denominator to $\sqrt{64+32}$. The pass matches supported split/concat graphs, retains rotary data and explicit scale, and does not claim arbitrary RoPE recognition. \textbf{Use case:} MLA-like decoupled rotary layouts.

\subsection*{3.4 Native-attention phase guard and decomposed fallback}
Online softmax may merge \emph{token} blocks after each score's full feature sum; it must not treat feature partial sums as extra tokens~\cite{online,flash}. In the seed-free 256-feature case, one 16-wide tile contributes scores $[\log9,0]$; the other 15 contribute $[0,0]$, with values $[1,0]$:
\[
\text{correct}=\frac9{9+1}=0.9,\qquad
\text{wrong feature-as-token merge}=\frac{9+15}{9+1+30}=0.6.
\]
The C++ configured-IR guard protects the feature reduction under the supported CPU schedule and resource constraints; a standard-contraction fallback offers another correct route. The later audit also finds an existing upstream early-decomposition route on CPU. This is a bounded correctness safeguard, not a new attention algorithm or a proven general GPU scheduling repair.

\subsection*{3.5 Bounded trig argument reduction: keep large-phase RoPE accurate}
\begin{lstlisting}
# Conceptual repaired path, for 32 < abs(x) <= 8192:
xd = f64(x)
r  = xd - floor(xd/(2*pi) + 0.5) * (2*pi)
y  = sin_f32(f32(r))             # analogously for cos
\end{lstlisting}
Higher-precision reduction limits cancellation before the existing FP32 trig evaluation. The pass retains the already-computed FP32 phase, supports static unencoded small kernels ($\leq4096$ elements) and falls back outside the guard. It is a bounded scheme, not full Payne--Hanek or an all-input accuracy proof~\cite{payne}. Unmodified model validation passed 5/9 natural cases; repair enabled 9/9 under the same tolerances. Both IREE timing variants use the repair and f64-preservation flags, isolating the cache transformation.
\newpage
\section{Evaluation protocol}
\smallnote{Standalone C++ MLIR 19 rewriting precedes IREE 3.13/LLVM 24 compilation (\code{89ece87b8985})~\cite{mlir}. GPU: Tesla T4, driver 580.82.07, FP32, TF32 off. SmolLM2-135M-Instruct~\cite{model}, checkpoint \code{12fd25f77366}: 30 layers, 9Q/3KV heads, width 64. Both IREE variants use O3, both GQA flags, identical split-argmax settings, trig repair and f64-preservation flags. Timing excludes compilation, prefill and upload: resident outputs, context synchronization, randomized order, 31 samples/repetition; median of nine repetition medians (three prompts, three repetitions). Nsight Compute 2025.3.1 app-range counters~\cite{ncu} use prompt 0, three independent processes and three calls/range, separately from timing. All 112 final counter exports were rechecked.}

\section{Winner factoring: pretrained two-candidate step}
{\small
\begin{tabularx}{\linewidth}{@{}r r r r r Y@{}}
\toprule
\textbf{Context} & \textbf{Base ms} & \textbf{Opt. ms} & \textbf{Speedup} & \textbf{Latency lower} & \textbf{Total DRAM change}\\
\midrule
2048 &43.264&31.815&\textbf{1.36$\times$}&\textbf{26.46\%}&$-3.91\%$\\
3072 &63.326&60.386&1.05$\times$&4.64\%&$+1.24\%$\\
4096 &82.409&77.858&1.06$\times$&5.52\%&$-0.27\%$\\
\bottomrule
\end{tabularx}}
\smallnote{This comparison isolates concat winner factoring within the repaired IREE pipeline. All natural and forced-winner cases pass, with exact cache prefixes and unchanged inputs. Worst absolute error: $5.05\!\times\!10^{-5}$ (atol$=5\!\times\!10^{-4}$, rtol$=3\!\times\!10^{-3}$). Ordinary decode has no matches. Writes fall but reads rise: full-model DRAM reduction is inconsistent. At context 2048, kernels increase 987$\to$1047/call and logical memcpy remains 188.8\,MB/call.}

\section{Per-transformation operator results}
\smallnote{Speedup $=$ baseline time / optimized time; below 1 means slower. CPU rows are separate synthetic studies, not full-model gains. Suite figures are geometric means. Earlier T4 projection timings include output downloads; final model timings leave outputs resident. Combined gains are not additive. Full per-case times, variant names and source paths: \code{report/study\_results/}.}
{\small
\begin{tabularx}{\linewidth}{@{}P{40mm}P{29mm}Y@{}}
\toprule
\textbf{Transformation} & \textbf{Measured effect} & \textbf{Workload / attribution}\hfill\cite{project}\\
\midrule
Original GQA head reuse &1.86$\times$ CPU&Default lowering only; no extra gain established over both upstream flags plus fusion.\\
Conditional read forwarding &2.04$\times$ CPU&Cache A1: 187.10$\to$91.74\,$\mu$s; same best-correct baseline.\\
Masked gather forwarding &2.52$\times$ CPU&Same A1: 187.10$\to$74.36\,$\mu$s; alternative to conditional reads.\\
Winner factoring alone &1.38$\times$ CPU&2607.34$\to$1885.37\,$\mu$s; with masked forwarding: 723.42\,$\mu$s, \textbf{3.60$\times$ combined}.\\
Adjacent-update coalescing &0.94--0.96$\times$ alone&4--16-token CPU chains. \textbf{2.17--2.34$\times$ combined} with masked forwarding; merging enables a later match.\\
Snapshot protection / COW &2.58$\times$ combined&109.01$\to$42.27\,$\mu$s CPU: protection \emph{plus} masked forwarding vs fastest correct baseline; not protection alone.\\
Key/value projection placement &3.93$\times$ CPU; 1.20$\times$ T4&16 cases per platform vs strong unabsorbed contractions. One decode ablation: key-only 1.92$\times$, value-only 1.95$\times$.\\
Partial RoPE + projection &11.40$\times$ CPU; 1.50$\times$ T4&Four cases per platform; combined pipeline, no isolated RoPE gain. MLA absorption is prior art.\\
Correct contraction fallback &10.63$\times$ CPU&One wide-head case: 131.11$\to$12.33\,$\mu$s vs correct guarded fused lowering; not a gain from the guard itself.\\
Gather normalization / native guard &Correctness enablers&No independent valid speedup established. Wrong stock outputs are excluded from performance baselines.\\
Stable trig reduction &5/9$\to$9/9 cases pass&Earlier matched-valid-case latency overhead: +0.07\% (A), +0.25\% (B); not a speed optimization.\\
Extra aggressive fusion flag &0.97$\times$ at 2048&Earlier repaired B run: 32.02$\to$33.16\,ms; does not consistently help.\\
\bottomrule
\end{tabularx}}

\section{Forwarding on T4: mechanism and negative control}
{\small
\begin{tabularx}{\linewidth}{@{}Y r r r P{36mm}@{}}
\toprule
\textbf{Track B fixture} & \textbf{Base ms} & \textbf{Opt. ms} & \textbf{Speedup} & \textbf{DRAM MB: base $\to$ opt.}\\
\midrule
Small sparse forks&0.288&0.246&1.17$\times$&5.831$\to$0.128 (97.80\% less)\\
Large sparse forks&1.777&1.230&1.44$\times$&178.897$\to$0.137 (99.92\% less)\\
Dense no-rewrite control&0.545&1.268&\textbf{0.43$\times$}&0.285$\to$0.275 (small / variable)\\
\bottomrule
\end{tabularx}}
\smallnote{Conditional forwarding, separate from the model factoring result. Track A large fixture: 1.14$\times$, 178.228$\to$0.132\,MB in a second T4 session. Synthetic checks use independent FP64 NumPy references (atol$=5\!\times\!10^{-5}$, rtol$=2\!\times\!10^{-3}$). The unexplained no-rewrite regression and uncontrolled clocks limit causal attribution.}

\textbf{Limitations.} One FP32 small model and GPU family; custom one-step branching, not complete serving throughput. Logical tensor elimination, physical allocation, DRAM bytes and latency are distinct claims. Incorrect ownership/numerical cases and irregular masked CUDA compile failures remain recorded. No all-GQA, mixed-precision or global-novelty claim is made.

\clearpage
\renewcommand{\refname}{References and attribution}
\smallnote{\textbf{Prior work.} GQA, MLA absorption and decoupled RoPE have identifiable originating papers below. Read-over-write has classical foundations. For selection factoring, adjacent-update merging and this exact compiler repair, no definitive first paper is established here. Cache-sharing papers are related systems precedents, not attribution for our precise MLIR patterns.}
\vspace{3pt}
\begingroup
\small
\setlength{\parskip}{1pt}
\begin{thebibliography}{99}
\setlength{\itemsep}{3pt}
\bibitem{gqa} J. Ainslie, J. Lee-Thorp, M. de Jong, Y. Zemlyanskiy, F. Lebr\'on and S. Sanghai.
\emph{GQA: Training Generalized Multi-Query Transformer Models from Multi-Head Checkpoints}. EMNLP, 2023, pp. 4895--4901. Originating GQA reference.
\url{https://aclanthology.org/2023.emnlp-main.298/}
\bibitem{mqa} N. Shazeer. \emph{Fast Transformer Decoding: One Write-Head is All You Need}. 2019. Introduces multi-query K/V sharing.
\url{https://arxiv.org/abs/1911.02150}
\bibitem{mccarthy} J. McCarthy. \emph{Towards a Mathematical Science of Computation}. Information Processing'62, 1962, pp. 21--28. Classical functional-array foundation; not a GQA optimization paper.
\url{https://www-formal.stanford.edu/jmc/towards/towards.html}
\bibitem{arrays} C. Tinelli / SMT-LIB. \emph{ArraysEx: Functional Arrays with Extensionality}. Formal select/store read-over-write axioms.
\url{https://smt-lib.org/theories-ArraysEx.shtml}
\bibitem{paged} W. Kwon et al. \emph{Efficient Memory Management for Large Language Model Serving with PagedAttention}. SOSP, 2023. KV sharing and copy-on-write precedent; not our compiler factoring pass.
\url{https://arxiv.org/abs/2309.06180}
\bibitem{forkkv} S. Wang, R. Ren and L. Gui. \emph{ForkKV: Scaling Multi-LoRA Agent Serving via Copy-on-Write Disaggregated KV Cache}. 2026. Related branching-cache system, not the first array-law reference.
\url{https://arxiv.org/abs/2604.06370}
\bibitem{deepseek} DeepSeek-AI et al. \emph{DeepSeek-V2: A Strong, Economical, and Efficient Mixture-of-Experts Language Model}. 2024. Secs.2.1.2--2.1.3 explicitly describe MLA weight absorption and decoupled RoPE.
\url{https://arxiv.org/abs/2405.04434}
\bibitem{rope} J. Su et al. \emph{RoFormer: Enhanced Transformer with Rotary Position Embedding}. First arXiv version, 2021. Originating RoPE reference.
\url{https://arxiv.org/abs/2104.09864}
\bibitem{payne} M. H. Payne and R. N. Hanek. \emph{Radian Reduction for Trigonometric Functions}. ACM SIGNUM Newsletter 18(1), 1983, pp. 19--24. Seminal argument-reduction reference; our bounded f64 scheme is not its full algorithm.
\url{https://doi.org/10.1145/1057600.1057602}
\bibitem{online} M. Milakov and N. Gimelshein. \emph{Online Normalizer Calculation for Softmax}. 2018. Online normalization background.
\url{https://arxiv.org/abs/1805.02867}
\bibitem{flash} T. Dao, D. Y. Fu, S. Ermon, A. Rudra and C. R\'e. \emph{FlashAttention: Fast and Memory-Efficient Exact Attention with IO-Awareness}. NeurIPS, 2022. Fused IO-aware attention background, not implemented anew in this project.
\url{https://arxiv.org/abs/2205.14135}
\bibitem{mlir} C. Lattner et al. \emph{MLIR: A Compiler Infrastructure for the End of Moore's Law}. 2020.
\url{https://arxiv.org/abs/2002.11054}
\bibitem{ireesource} IREE contributors. \emph{GeneralizeLinalgNamedOps; BubbleUpExpandShapes; compilation pipelines}. Source pinned to compiler commit \code{89ece87b89852100c0a4ac3d515e14842d090967}. Implementation evidence for existing transformations, not a GQA-specific paper.
Pinned source: \href{https://github.com/iree-org/iree/blob/89ece87b89852100c0a4ac3d515e14842d090967/compiler/src/iree/compiler/GlobalOptimization/GeneralizeLinalgNamedOps.cpp}{GeneralizeLinalgNamedOps.cpp} and \href{https://github.com/iree-org/iree/blob/89ece87b89852100c0a4ac3d515e14842d090967/compiler/src/iree/compiler/DispatchCreation/BubbleUpExpandShapes.cpp}{BubbleUpExpandShapes.cpp}.
\bibitem{ireedocs} IREE contributors. \emph{MLIR pass reference: GlobalOptimization and DispatchCreation}. Named-op generalization, reshape propagation through reductions and elementwise fusion. The observed behavior is pinned by the source/IR evidence, since online documentation can change.
\url{https://iree.dev/reference/mlir-passes/GlobalOptimization/}
\url{https://iree.dev/reference/mlir-passes/DispatchCreation/}
\bibitem{model} Hugging Face. \emph{SmolLM2-135M-Instruct model card and checkpoint}. Pinned revision in the evaluation protocol.
\url{https://huggingface.co/HuggingFaceTB/SmolLM2-135M-Instruct}
\bibitem{ncu} NVIDIA. \emph{Nsight Compute Profiling Guide}. Application-range replay, metric coverage and cache/clock controls.
\url{https://docs.nvidia.com/nsight-compute/ProfilingGuide/}
\bibitem{project} CS6886 project. \emph{Attention and KV-cache compiler optimizations: implementation, raw results and reproduction}. Experiment-source snapshot \code{71cc40ae53c5e60f2a690e67ffbc64e85969e64f}.
\url{https://github.com/lazarsalt1441/cs6886_project}
\end{thebibliography}
\endgroup
\end{document}
